\section{Discussion}
\label{sec:discussion}

A single stochastic attribution is conditionally adequate for the explicit numerical reading tested here. At the highest budgets it always met the broad magnitude/two-of-three signed-selection criterion in the observed panels; lower budgets, particular fitted predictors and stricter feature-list requirements produced failures. This supports budgeting and checking explanations against the reading they must preserve. It does not support trusting every single run because SHAP satisfies an additive identity, because another run looks similar, or because the panel average is high.

An assessment should specify the probability game and feature units, audit a feasible reference for that game, measure both vector and signed-selection error, and show the distributions across inputs. Work and CPU time make its resource requirement checkable. The changed-background arm further shows why accurate estimation for one chosen empirical game cannot substitute for justifying that game. If the application requires all three listed signed features to be recovered, our broad event is insufficient. If it requires a population or per-prediction reliability certificate, neither the finite observed all-success panels nor thirty seeds supply one. The measured failure cases and inactive low-signal sensitivity also limit any attempt to transfer these numerical tolerances directly to new inputs or tasks.

The contribution is this executed operational assessment and its reproducible evidence, not a new estimator or a universal tolerance. The comparison extends established estimator-error and convergence work to a fully enumerated, fixed-game tabular panel with explicit one-run criteria, matched input uncertainty, separate repeatability, actual computation accounting and retained negative controls. Numerical reliability remains only one requirement for using an attribution as a reason for a consequential prediction.

\paragraph{Broader impacts.}
Showing when a displayed attribution changes can help users avoid overconfident numerical explanations and make computation requirements transparent. Conversely, an apparently stable feature list can lend unjustified authority to a biased credit or health prediction; even perfect agreement with our empirical game would not make that decision fair or appropriate. We use existing public benchmarks without contacting or recruiting participants, do not redistribute raw person records or release fitted decision services, and retain source-use qualifications and demographic/survey limitations. The study's modest CPU use is reported rather than converted into an unsupported emissions estimate.
